\documentclass[10pt,a4paper]{amsart}
\usepackage[margin=19mm]{geometry}
\usepackage{amsmath,amssymb,mathtools}
\usepackage[scr=rsfs,cal=euler]{mathalfa}
\usepackage{xfrac}
\usepackage[unicode,hidelinks,bookmarks=false]{hyperref}
\newcommand{\ie}{i.e.\ }
\newcommand{\cf}{cf.\ }
\newcommand{\resp}{respectively\ }
\newcommand{\Subsection}{\subsection}
\providecommand{\nobreakdash}{\nobreak\hbox{-}}
\setlength{\emergencystretch}{2em}
\multlinegap0pt
\usepackage{mdframed}
\usepackage{mdframed}
\usepackage{mdframed}
\usepackage{mdframed}
\usepackage{mathtools}
\usepackage{amssymb}
\usepackage{tikz}
\usepackage{xcolor}
\usepackage{algorithm}\usepackage{algpseudocode}
\usepackage{mdframed}
\usepackage[shortlabels]{enumitem}
\newtheorem{lemma}{Lemma}
\newtheorem{definition}{Definition}
\usepackage{cleveref}
\crefname{lemma}{Lemma}{Lemmas}
\crefname{definition}{Definition}{Definitions}
\newcommand{\E}{\mathbb E}
\newcommand{\T}{\mathcal E}
\title{Wasserstein Exponential Smoothing for Distributional Time Series Forecasting}
\author{Takuo Matsubara, Peiwen Jiang, Minh-Ngoc Tran, Wilson Ye Chen}
\date{Source: arXiv:2606.05560v2 (CC BY 4.0)}
\begin{document}
\maketitle

\noindent\emph{Source and licence.} Wasserstein Exponential Smoothing for Distributional Time Series Forecasting. Version arXiv:2606.05560v2 (CC BY 4.0), distributed by arXiv under the Creative Commons Attribution 4.0 International licence, http://creativecommons.org/licenses/by/4.0/, as recorded in the arXiv licence metadata for this exact version and shown on the abstract page https://arxiv.org/abs/2606.05560v2. Full licence notice: \texttt{licenses/arxiv-2606.05560-CC-BY-4.0.txt}.\par\smallskip

\noindent\textit{Editorial scope.} Counted unit: the article's lemma prop:model\_quantile\_decomposition (source0.tex lines 1075--1080). The article's complete original proof of it is delivered with it (source0.tex lines 1367--1420, one \texttt{proof} environment of 54 lines, which the article places away from the statement). No mathematics is written, repaired or re-derived: the statement and the proof are the article's own lines, in the article's own notation, and no retained line is substituted or re-flowed.  Retained uncounted as the source-local context the proof needs: lines 1000--1000; lines 1032--1042; lines 1064--1071; lines 1104--1115.  Nothing was omitted from any retained span, so the package carries no omission record.  No retained line contains a citation, so the wrapper carries no bibliography and no bibliography entry.  No retained line contains a figure environment, an image-inclusion command or drawing code, so no image is delivered and the package declares no asset dependency.  Editorial formatting only, in addition: an AMS \texttt{amsart} preamble that restates the article's own declarations and macros used by the retained lines, namely the environments \texttt{lemma}, \texttt{definition} on separate counters, together with the standard packages those lines need.  The retained environments print in this document as Lemma 1, Definition 1 in order of appearance, whereas the article prints the same items with the numbering of its own full document.  The complete native source of the article as distributed in the arXiv e-print for this exact version is bundled unchanged as \texttt{sources/202125/source0.tex}.
\section{Preliminaries for Main Proofs} \label{apx:preparation}

\begin{lemma} \label{lem:residual_property}
	For each $t \ge 1$, the quantile residual $F_{t}$ is a random element defined by the sequence of random transport maps $\T_1, \dots, \T_{t}$.
	For all $t \ge 1$, the quantile residual $F_{t}$ satisfies:
	\begin{enumerate}
		\item $\E\big[ F_{t}(q) \mid \T_1, \dots, \T_{t-1} \big] = 0$ at every $q \in (0, 1)$;
		\item $\E\big[ \langle F_{t}, F_{s} \rangle_2 \mid \T_1, \dots, \T_{t-1} \big] = 0$ for any time $s$ such that $1 \le s < t$;
		\item $\E\big[ \| F_{t} \|_2^2 \mid \T_1, \dots, \T_{t-1} \big] = \sigma_{t}^2$;
		\item $\E\big[ \| F_{t} \|_2^4 \mid \T_1, \dots, \T_{t-1} \big] \le \kappa$ .
	\end{enumerate}
	Furthermore, these imply: $\E[ F_{t}(q) ] = 0$, $\E[ \langle F_{t}, F_{s} \rangle_2 ] = 0$, $\E[ \| F_{t} \|_2^2 ] = \sigma_{t}^2$, and $\E[ \| F_{t} \|_2^4 ] \le \kappa$.
\end{lemma}

\begin{definition}[Quantile Residual Cumulative] \label{def:residual_cumulative}
	For each $t \ge 1$ and $\theta \in (0, 1)$, let $G_t^\theta$ be a random element in $L^2$ defined by the sequence of quantile residuals $F_1, \dots, F_{t}$, given by
	\begin{align}
		G_t^\theta(q) := \sum_{i=1}^{t} (1 - \theta)^{t-i} F_i(q) ,
	\end{align}
    where we set $G_0^\theta(q) := 0$ for $t = 0$.
	We call $G_t^\theta$ the quantile residual cumulative at time $t$.
\end{definition}

\begin{lemma} \label{lem:g_moment}
	At each time $t \ge 1$ and $\theta \in (0, 1)$, the quantile residual cumulative $G_t^\theta$ is a random element defined by the sequence of random transport maps $\T_1, \dots, \T_t$. 
    Define $m_t^\theta := \sum_{i=1}^{t} (1 - \theta)^{2(t-i)} \sigma_i^2$, setting $m_0^\theta := 0$ for $t = 0$.
    For any $0 \le s < t$, there exists a time-independent constant $c^\theta$, such that:
	\begin{enumerate}
        \item $\E[ G_t^\theta(q) ] = 0$;
        \item $\E[ \langle F_t, G_s^\theta \rangle ] = 0$;
		\item $\E[ \| G_t^\theta \|_2^2 ] = m_t^\theta \le \frac{\sigma^2}{\theta (2 - \theta)}$;
		\item $\E[ ( \| G_t^\theta \|_2^2 - m_t^\theta )^2 ] \le c^\theta$;
		\item $\E[ ( \| G_t^\theta \|_2^2 - m_t^\theta ) ( \| G_s^\theta \|_2^2 - m_s^\theta ) ] \le (1 - \theta)^{2(t-s)} c^\theta$ .
	\end{enumerate}
\end{lemma}

\begin{lemma} \label{prop:model_quantile_decomposition}    
	For any $t \ge 0$ and $\theta \in (0, 1)$, we have for all $q \in (0, 1)$:
	\begin{align}
		U_{t}^\theta(q) - V_{t+1}(q) = (1 - \theta)^{t} \big( U_0^\theta(q) - U_0(q) \big) + ( \theta - \theta_* ) G_{t}^\theta(q) - F_{t+1}(q) .
	\end{align}
\end{lemma}

\begin{proof}
Let $\Delta_0 := U_0^\theta - U_0$, noting that $\Delta_0$ is deterministic and independent of $\theta$, since $U_0^\theta$ is the initial condition of the WES filter.
\Cref{prop:model_quantile_decomposition} (see \Cref{apx:preparation}) establishes that
\begin{align}
    U_{t}^\theta - V_{t+1} = (1 - \theta)^t \Delta_0 + ( \theta - \theta_* ) G_t^\theta - F_{t+1} .
\end{align}
Consequently, the autocovariance of interest $\E[ \langle U_{t}^\theta - V_{t+1},U_{s}^\theta - V_{s+1} \rangle ]$ expands to
\begin{align}
    (*) & := \E[ \langle \, (1 - \theta)^t \Delta_0 + ( \theta - \theta_* ) G_t^\theta - F_{t+1} \, , \, (1 - \theta)^s \Delta_0 + ( \theta - \theta_* ) G_s^\theta - F_{s+1} \, \rangle ] \\
    & = \E\left[ \langle (1-\theta)^t \Delta_0, (1-\theta)^{s} \Delta_0 \rangle \right] + \E\left[ \langle (1-\theta)^t \Delta_0, (\theta - \theta_*) G_s^\theta \rangle \right] + \E\left[ \langle (1-\theta)^t \Delta_0, F_{s+1} \rangle \right] \\
    & \hspace{25pt} + \E\left[ \langle (\theta - \theta_*) G_t^\theta, (1-\theta)^{s} \Delta_0 \rangle \right] + \E\left[ \langle (\theta - \theta_*) G_t^\theta, (\theta - \theta_*) G_s^\theta \rangle \right] + \E\left[ \langle (\theta - \theta_*) G_t^\theta, F_{s+1} \rangle \right] \\
    & \hspace{50pt} + \E\left[ \langle F_{t+1}, (1-\theta)^{s} \Delta_0 \rangle \right] + \E\left[ \langle F_{t+1}, (\theta - \theta_*) G_s^\theta \rangle \right] + \E\left[ \langle F_{t+1}, F_{s+1} \rangle \right] \\
    & = (1-\theta)^{t+s}\|\Delta_0\|^2 + (1-\theta)^t (\theta - \theta_*) \left\langle \Delta_0, \E\big[ G_s^\theta \big] \right\rangle + (1-\theta)^t \left\langle \Delta_0, \E\big[ F_{s+1} \big] \right\rangle \\
    & \hspace{25pt} + (1-\theta)^s (\theta - \theta_*) \left\langle \E\big[ G_t^\theta \big], \Delta_0 \right\rangle + (\theta - \theta_*)^2 \E\left[ \langle G_t^\theta, G_s^\theta \rangle \right] + (\theta - \theta_*) \E\left[ \langle G_t^\theta, F_{s+1} \rangle \right] \\
    & \hspace{50pt} + (1-\theta)^{s} \left\langle \E\big[ F_{t+1} \big], \Delta_0 \right \rangle + (\theta - \theta_*) \E\left[ \langle F_{t+1}, G_s^\theta \rangle \right] + \E\left[ \langle F_{t+1}, F_{s+1} \rangle \right] .
\end{align}
By \Cref{lem:residual_property} and \Cref{lem:g_moment}, we have $\E[ G_s^\theta ] = 0$, $\E[ F_{s+1} ] = 0$, $\E[ G_t^\theta ] = 0$, $\E[ F_{t+1} ] = 0$, $\E[ \langle F_{t+1}, G_s^\theta \rangle ] = 0$, and $\E[ \langle F_{t+1}, F_{s+1} \rangle ] = 0$.
Substituting these identities simplifies $(*)$ to
\begin{align}
    (*) & = (1-\theta)^{t+s}\|\Delta_0\|^2 + (\theta - \theta_*)^2 \underbrace{ \E\left[ \langle G_t^\theta, G_s^\theta \rangle \right] }_{ (*_1) } + (\theta - \theta_*) \underbrace{ \E\left[ \langle G_t^\theta, F_{s+1} \rangle \right] }_{ (*_2) } . \label{eq:autocovariance_decomp}
\end{align}
We evaluate the terms $(*_1)$ and $(*_2)$ in the remainder.

\vspace{5pt}
\noindent
\textbf{Term $(*_1)$:}
Recall the definition of $G_t^\theta$ in \Cref{def:residual_cumulative}.
The following decomposition holds:
\begin{align}
    G_t^\theta = \sum_{i=1}^{s} (1-\theta)^{t-i} F_i + \sum_{i=s+1}^{t} (1-\theta)^{t-i}F_i = (1-\theta)^{t-s} \underbrace{ \sum_{i=1}^s (1-\theta)^{s-i} F_i }_{ = G_s^\theta } + \sum_{i=s+1}^t (1-\theta)^{t-i} F_i
\end{align}
It then follows from this decomposition and \Cref{lem:g_moment} that
\begin{align}
    (*_1) = \E\left[ \langle G_t^\theta, G_s^\theta \rangle \right] = (1 - \theta)^{t-s} \E\big[ \| G_s^\theta \|^2 \big] + \sum_{i=s+1}^t (1 - \theta)^{t-i} \E\big[ \langle G_s^\theta, F_i \rangle \big] = (1-\theta)^{t-s} m_s^\theta ,
\end{align}
where $m_s^\theta = \sum_{i=1}^{t} (1 - \theta)^{2(t-i)} \sigma_i^2$ as defined in \Cref{lem:g_moment}.

\vspace{5pt}
\noindent
\textbf{Term $(*_2)$:}
Returning to the definition of $G_t^\theta$, we obtain
\begin{align}
    (*_2) = \E\left[ \langle G_t^\theta, F_{s+1} \rangle \right] = \sum_{i=1}^{t} (1-\theta)^{t-i} \E\left[ \langle F_i, F_{s+1} \rangle \right] .
\end{align}
\Cref{lem:residual_property} guarantees that $\E\left[ \langle F_t, F_{s+1} \rangle \right] = 0$ for all $i \ne s + 1$, meaning that only $i = s + 1$ remains:
\begin{align}
    (*_2) = (1 - \theta)^{t-s-1} \E\left[ \| F_{s+1} \|_2^2 \right] = (1 - \theta)^{t-s-1} \sigma_{s+1}^2 .
\end{align}
where the second equality follows from $\E[ \| F_{s+1} \|_2^2 ] = \sigma_{s+1}^2$ by \Cref{lem:residual_property}.

\vspace{5pt}
\noindent
Substituting the equalities of $(*_1)$ and $(*_2)$ into \eqref{eq:autocovariance_decomp} completes the proof.
\end{proof}

\end{document}
